\documentclass[11pt,letterpaper]{article}
\usepackage[T1]{fontenc}
\usepackage[utf8]{inputenc}
\usepackage{lmodern,microtype}
\usepackage[margin=0.79in,headheight=15pt]{geometry}
\usepackage{parskip,enumitem,longtable,tabularx,booktabs,array,xcolor,fancyhdr,hyperref}
\definecolor{navy}{HTML}{17314C}
\definecolor{teal}{HTML}{087F78}
\hypersetup{colorlinks=true,linkcolor=navy,urlcolor=teal}
\setlength{\emergencystretch}{3em}
\setlist[itemize]{leftmargin=*,nosep,topsep=5pt}
\setlist[enumerate]{leftmargin=*,itemsep=3pt}
\renewcommand{\arraystretch}{1.22}
\pagestyle{fancy}\fancyhf{}
\fancyhead[L]{\small Enterprise Data Modeling}\fancyhead[R]{\small Documentation Suite}
\fancyfoot[C]{\thepage}
\newcommand{\srcbook}{\textit{The Data Model Resource Book}, revised edition, Volume 2}
\newcommand{\ResponseLine}{\par\vspace{1.05em}\noindent\rule{\linewidth}{0.4pt}\par}
\newcommand{\ResponseArea}{\ResponseLine\ResponseLine\ResponseLine}
\begin{document}
\begin{center}
{\LARGE\bfseries\color{navy} Classification, Lineage and Retention Controls}\par
\vspace{0.25em}{\large Supplementary governance reference}\par
\vspace{0.3em}{\small Original engineering study material \textbullet{} 2026 edition}
\end{center}
\vspace{0.8em}
\section{Scope}
Supplementary enterprise governance practice applied to universal and industry models. The supplied 2001 source is not treated as a modern regulatory or security policy. Consult the organization's legal and compliance specialists for binding requirements.


\section{Operating controls}
\begin{tabularx}{\linewidth}{@{}p{0.24\linewidth} p{0.72\linewidth}@{}}
\toprule
\textbf{Control domain} & \textbf{Required practice} \\ 
\midrule
Classification & Assign categories according to organizational policy and domain law; never assume a data element is public because it appears in a test sample. \\ 
Minimization & Store what is needed for a defined purpose; prefer derived or tokenized values where feasible. \\ 
Lineage & Capture source system, extraction/transform version, creation time, responsible owner and recipient. \\ 
Access & Apply least privilege using approved identity and authorization systems; test row and column boundaries where appropriate. \\ 
Retention & Record trigger event, minimum/maximum retention, legal hold, archive and deletion evidence; seek legal review for regulated records. \\ 
Operational safety & Sanitize logs, backups and test fixtures; protect secrets and identifiers from repository leaks. \\ 
\bottomrule
\end{tabularx}

\section{Required evidence register}
\begin{itemize}
\item Entity/data product owner, steward, accountable system and current version.
\item Sensitivity, permitted use, geographic/jurisdictional restrictions and retention trigger.
\item Required data quality assertions and current pass/fail evidence.
\item Data lineage upstream and downstream, including intersystem ID translation.
\item Approval references for breaking changes, access grants and policy exceptions.
\item Verification dates and authoritative links, with no passwords or sensitive data in the repository.
\end{itemize}

\section{Review cadence}
Review on design introduction, material schema or semantic change, incident, and scheduled policy attestation. The review records what was checked, the evidence date, unresolved risks, remediation owner and next checkpoint.


\vfill
\hrule
{\footnotesize\textbf{Provenance.} Supplementary engineering practice: original content; the cited book provides domain-modeling context.}
\end{document}
